% !TeX program = xelatex
\documentclass[a4paper,11pt]{article}
\usepackage[margin=27mm,headheight=14pt]{geometry}
\usepackage{amsmath,amssymb,amsthm,mathtools,bm}
\usepackage{booktabs,array,longtable}
\usepackage{xcolor,xurl}
\definecolor{linkblue}{RGB}{25,64,117}
\usepackage[colorlinks=true,linkcolor=linkblue,citecolor=linkblue,urlcolor=linkblue,bookmarksnumbered=true,
 pdftitle={A 17/4 Hausdorff bound for six-dimensional Kakeya sets},
 pdfsubject={A derivation from a quantitative four-dimensional input}]{hyperref}
\usepackage{fancyhdr}
\pagestyle{fancy}\fancyhf{}
\fancyhead[L]{\small A $17/4$ bound for six-dimensional Kakeya sets}
\fancyhead[R]{\small Research manuscript}
\fancyfoot[C]{\thepage}\renewcommand{\headrulewidth}{0.3pt}
\setlength{\parskip}{0.25em}\setlength{\emergencystretch}{2em}
\numberwithin{equation}{section}
\theoremstyle{plain}
\newtheorem{theorem}{Theorem}[section]
\newtheorem{proposition}[theorem]{Proposition}
\newtheorem{lemma}[theorem]{Lemma}
\newtheorem{corollary}[theorem]{Corollary}
\theoremstyle{definition}
\newtheorem{definition}[theorem]{Definition}
\newtheorem{inputthm}[theorem]{Cited input}
\theoremstyle{remark}\newtheorem{remark}[theorem]{Remark}
\newcommand{\R}{\mathbb R}\newcommand{\T}{\mathcal T}\newcommand{\Kest}{\mathsf K}
\newcommand{\HH}{\operatorname{dim}_{\mathrm H}}\newcommand{\ang}{\operatorname{ang}}
\newcommand{\card}[1]{\#\!\left(#1\right)}
\newcommand{\norm}[1]{\left\lVert #1\right\rVert}
\newcommand{\Ldelta}{L_\delta}\DeclareMathOperator{\mass}{mass}
\title{A $17/4$ Hausdorff bound for six-dimensional Kakeya sets\\from a quantitative four-dimensional input}
\author{}\date{October 26, 2026}
\begin{document}
\maketitle
\begin{abstract}
Starting from an explicitly stated four-dimensional finite direction-mass estimate, we derive a Hausdorff dimension lower bound of $17/4$ for every set $K\subset\R^6$ containing a unit line segment in every direction. No measurability, compactness, boundedness, or measurable choice of segments is assumed. The finite step projects a direction family in an arbitrary fixed ambient space to $\R^4$. Gaussian small-ball estimates and a collision-graph independent set simultaneously preserve separated directions, nonshort projections, and marks belonging to their original lines. For average shading density $\lambda$, angular upper-bound constant $A$, and angular exponent $s\ge3$, the resulting estimate retains the factor $\lambda^2\card{\T}/A$ and an endpoint logarithmic loss. This provides the unsaturated intermediate estimates required by the Katz--Tao recurrence. An outer-measure argument then transfers the finite estimates to arbitrary sets. We explicitly cite the four-dimensional proof-level finite interface and the Katz--Tao theorem; a four-dimensional set-dimension statement alone is not substituted for the quantitative input. The exponent exceeds the explicitly checked six-dimensional baseline $7-2\sqrt2$ by approximately $0.078427$.
\end{abstract}
\noindent\textbf{Keywords:} Kakeya sets; Hausdorff dimension; random projection; finite-scale estimates; Katz--Tao recurrence.

\medskip
\noindent\textbf{Status and dependencies.}
This is a manuscript presentation of the existing research record DIM-001. It has undergone internal cross-review, but has not been externally peer reviewed. Appendix~\ref{app:inputs} identifies the four-dimensional finite input, its advanced dependencies, and the archived proofs. We do not independently reprove the complete four-dimensional manuscript or the Katz--Tao recurrence, assert absolute priority or current optimality, or resolve the full-dimensional conjecture in general dimensions. The literature-comparison cutoff is October 19, 2026; the preparation date does not indicate a new frontier search.

\section{Introduction and main result}
\label{sec:introduction}

A Kakeya set in real Euclidean space is a set containing a unit line segment in every unoriented direction.
The Kakeya conjecture predicts that every such set in $\R^n$ has Hausdorff dimension $n$.
Throughout this paper we consider ordinary straight-line Kakeya sets in real Euclidean space.
We distinguish this problem from its finite-field analogue, Minkowski dimension bounds,
restricted-direction results, and results for classes of sets satisfying additional structural assumptions.

The four-dimensional manuscript\cite{OAI4} claims that every four-dimensional Kakeya set has full Hausdorff dimension.
For a transfer to higher dimensions, this set-dimension statement alone is insufficient:
the recurrence requires a quantitative input applicable to arbitrary finite line families and their original shadings
in every fixed ambient dimension, with constants allowed to depend on that dimension.
We use a finite-scale consequence of that proof, together with its sparse-direction amplification,
stated explicitly as Input~\ref{input:D4}.
Its proof dependencies and archival references are described in Appendix~\ref{app:inputs}.
The existing finite proof and outcome report are recorded in\cite{ArchiveFSD,ArchiveDIM}, respectively.

\begin{theorem}[Six-dimensional lower bound]\label{thm:main}
Assume the four-dimensional finite input~\ref{input:D4} and the Katz--Tao recurrence theorem~\ref{input:KT}.
If an arbitrary set $K\subset\R^6$ contains a unit line segment in every Euclidean direction, then
\begin{equation}\label{eq:main}
  \HH K\ge\frac{17}{4}=4.25.
\end{equation}
The set $K$ need not be measurable, compact, or bounded, and no measurable direction-to-line selection is required.
\end{theorem}

The proof has three steps. First, random linear projection sends a finite line family in any fixed $\R^M$ to $\R^4$.
Projected directions may collide, and projected segments may become short; both losses must be accounted for before invoking the four-dimensional input.
Second, an unsaturated estimate retaining the arbitrary original number of lines yields the strictly subcritical seed
$\Kest(n,4-\zeta)$ for every fixed real parameter $n\ge4$.
For $0<\zeta<1$, the Katz--Tao theorem gives
\begin{equation}\label{eq:roadmap-recurrence}
  \Kest(6,4-\zeta),\quad\Kest(5-\zeta,4-\zeta)
  \quad\Longrightarrow\quad
  \Kest\!\left(6,\frac{17}{4}-\frac{\zeta}{4}\right).
\end{equation}
Third, covers at a common scale and outer measure on the direction space yield a Hausdorff dimension bound for arbitrary sets; we then let $\zeta\downarrow0$.

The classical six-dimensional bound of Katz--Tao\cite{KT} is $7-2\sqrt2$.
The corresponding Euclidean comparisons of Hickman--Rogers--Zhang\cite[Figure 5]{HRZ}
and Gao--Liu--Xi\cite[Table 1]{GLX} also explicitly list this value.
The difference is
\begin{equation}\label{eq:comparison}
 \frac{17}{4}-(7-2\sqrt2)
   =2\sqrt2-\frac{11}{4}
   \approx0.078427125>0.
\end{equation}
This comparison concerns only the explicitly checked ordinary six-dimensional Hausdorff baseline;
it is not a certification of optimality across the entire current literature.
The contribution presented here is a higher-dimensional consequence of the four-dimensional finite input
and an exposition of the finite interfaces needed for this consequence. The recurrence itself is due to Katz--Tao.

\paragraph{Organization.}
Section~\ref{sec:preliminaries} fixes the discrete conventions and states the two named inputs.
Section~\ref{sec:projection} proves the finite random projection estimate while retaining the original shadings.
Section~\ref{sec:seeds} gives the unsaturated estimate and the actual parameters used in the recurrence.
Section~\ref{sec:hausdorff} establishes the Hausdorff conclusion for arbitrary sets and completes the proof of Theorem~\ref{thm:main}.
Section~\ref{sec:scope} explains the scope; the appendix records the source of the four-dimensional input and its sparse reduction.

\section{Finite-scale conventions and cited inputs}\label{sec:preliminaries}

\subsection{Directions, line families, and shadings}
Every ambient dimension $M$ is a fixed integer. Implicit constants may depend on $M$ and on a fixed spatial bound $R$,
but $M$ must not vary with the discretization scale.
Let $0<\delta<1/2$, and write
\[
  \Ldelta=1+\log(1/\delta).
\]
For unoriented unit directions $u,v$, use the projective angle
$\ang(u,v)=\arccos|\langle u,v\rangle|\in[0,\pi/2]$.
A set of directions is $\delta$-separated if every two distinct directions have projective angle at least $\delta$.

Let $\T$ be a finite labelled family of unit-scale line segments in a fixed ball $B(0,R)\subset\R^M$,
or of their neighbourhoods of width a fixed multiple of $\delta$.
A shading $Y$ assigns to each original line $T$ a collection of $\delta$-cells intersecting it or lying within distance $O(\delta)$ of it.
Write
\[
 m=\card{\T},\qquad E=\bigcup_{T\in\T}Y(T),\qquad
 S=\sum_{T\in\T}\card{Y(T)}.
\]
Here $\card{E}$ denotes the number of cells, not their Lebesgue volume.
Each original line has a fixed upper bound on its length, so
\begin{equation}\label{eq:single-line-upper}
  \card{Y(T)}\le C_{M,R}\delta^{-1}.
\end{equation}
Moving each cell to the nearest parameter point on its original line and removing multiplicities in parameter intervals of length $\delta$
loses only a fixed factor. We may therefore count separated marks on the original lines instead.
Throughout, cells and marks are interchanged only up to this fixed factor, without changing the small-power exponents.

For a real number $s\ge3$ and $A\ge1$, the required directional hypothesis is
\begin{equation}\label{eq:angular-cap}
 \card{\{T\in\T:\ang(v_T,u)\le\theta\}}
 \le A\left(\frac\theta\delta\right)^s,
 \qquad \delta\le\theta\le1,
\end{equation}
for every direction centre $u$.
This is only an upper bound; no pointwise Ahlfors--David lower bound, convex Wolff condition, or sticky hypothesis is imposed.
An average shading density $\lambda$ means
\begin{equation}\label{eq:mean-density}
 S\ge\lambda m\delta^{-1},\qquad 0<\lambda\le1.
\end{equation}
An average lower bound cannot directly replace a lower bound for each line.

\subsection{The four-dimensional finite directional-mass input}

\begin{inputthm}[Four-dimensional finite directional-mass estimate, archive DIR-003]\label{input:D4}
Fix a bounded four-dimensional graph chart. For every $\varepsilon_D>0$, there exist
$\eta_D>0$, $\delta_D>0$, and $C_D$ such that the following holds whenever $0<\delta<\delta_D$.
Let
\[
 \ell_i(t)=(t,b_i+t v_i),\qquad 0\le t\le1,\qquad
 |b_i|+|v_i|\le R_D,
\]
be a finite labelled line family, and let $J_i$ be its respective marks on the $\delta$ time grid, with
$\card{J_i}\ge\delta^{-1+\eta_D}$.
Given arbitrary positive weights $\omega_i$, suppose that every half-open slope cell $Q$ of side length $\delta$
satisfies $\sum_{v_i\in Q}\omega_i\le\alpha$.
Then the union $E_D$ of the cells visited by all the marks satisfies
\begin{equation}\label{eq:D4}
 \sum_i\omega_i\card{J_i}
 \le C_D\alpha\delta^{-\varepsilon_D}\card{E_D}.
\end{equation}
Repeated labels and arbitrary ranges of positive weights are allowed; the constants do not depend on the smallest weight or the number of labels.
Fixed-factor enlargements of the grid and visited cells are allowed, with corresponding changes in the constants.
\end{inputthm}

This input follows from a finite-scale consequence of the four-dimensional proof and a directional amplification;
see\cite{ArchiveSource,ArchiveDIR}.
It is not deduced directly from the set-dimension statement of\cite[Theorem 1.1]{OAI4} alone.
We use\eqref{eq:D4} with unit weights and separated output directions.
In particular, the input retains the order $\forall\varepsilon_D\,\exists\eta_D$;
no linear lower bound for $\eta_D$ in terms of $\varepsilon_D$ is assumed.

\subsection{Katz--Tao notation for estimates}

Set $N=\delta^{-1}$. We use the discrete-line, shading, and small-power-loss conventions of\cite[Definitions 5.1, 5.2, 6.1]{KT}.
For real $n,d>0$, $\Kest(n,d)$ means that, in every fixed ambient space $\R^M$,
a saturated $(n-1)$-collection of lines with saturated shading visits at least $N^{d-o(1)}$ cells.
The directional family satisfies the angular upper bound $N^{o(1)}(N\theta)^{n-1}$
for all centres and all $N^{-1}\le\theta\le1$, and has size $N^{n-1+o(1)}$.
Saturation of the shading means that its total mass is $N^{1-o(1)}\card{\T}$; saturation on each line is not assumed in advance.
The parameter $n$ is not the ambient dimension $M$ and need not be an integer.

Here $o(1)$ denotes only the arbitrarily small power losses allowed in the cited paper.
In an application, one first fixes the target output loss, then chooses a sufficiently small input tolerance, and finally takes $N$ sufficiently large.
Section~\ref{sec:seeds} gives an explicit $\forall\varepsilon\,\exists\eta$ sufficient form for this use.
We invoke only the strictly subcritical $\Kest(n,4-\zeta)$, without asserting a uniform linear modulus relating the critical input and output losses.

\begin{inputthm}[Katz--Tao recurrence, Theorem 6.2 of\cite{KT}]\label{input:KT}
Let $0<d<n$. If $\Kest(n,d)$ and $\Kest(d+1,d')$ hold, then
\begin{equation}\label{eq:KT}
  \Kest\!\left(n,\frac{2n+1+d'}4\right)
\end{equation}
holds. The input allows real parameters and every fixed ambient space.
\end{inputthm}

The second application in this recurrence may take place in $\R^{M+1}$ and requires the unsaturated version in equation (40) of the cited paper:
\begin{equation}\label{eq:KT-unsaturated}
  \card{E}\ge
  \frac{\card{\T}}{N^{n-1}}N^{d-o(1)}.
\end{equation}
Thus a saturated conclusion in just one fixed ambient space is insufficient.
Below we retain the original number-of-lines factor directly, establishing this interface;
the full proof of the recurrence remains a cited input from\cite{KT}.

\section{Finite random projection and retention of the original shadings}\label{sec:projection}

\begin{lemma}[Directional collisions and nonshort projections]\label{lem:projection}
Fix $M\ge4$, $s\ge3$, and a constant $a>0$.
If $m$ original $\delta$-separated directions satisfy\eqref{eq:angular-cap}, then there exist
$P:\R^M\to\R^4$ with $\norm P\le C_M$ and a subfamily $\T_*$ of the original labels such that
\begin{align}
 |Pv_T|&\ge c_M &&(T\in\T_*),\label{eq:nonshort}\\
 \ang(Pv_T,Pv_{T'})&>a\delta &&(T\ne T'\in\T_*),\label{eq:projected-separated}\\
 \card{\T_*}&\ge
 \frac{c_{M,s,a}m}{A\delta^{3-s}\Ldelta}.\label{eq:projection-count}
\end{align}
\end{lemma}

\begin{proof}
Let the entries of $P$ be independent standard real Gaussian variables.
Choose $B_M$ sufficiently large that $\Pr\{\norm P>B_M\}<1/32$.
For a fixed unit vector $u$, $Pu$ is a standard four-dimensional Gaussian vector.
We can therefore choose $c_M>0$ sufficiently small that $\Pr\{|Pu|<c_M\}<1/128$.

If $u,v$ have projective angle $\theta$, then, after changing the sign of $v$ if necessary, we can write
\[
 Pv=(\cos\theta)Pu+(\sin\theta)H,
\]
where $H$ is a standard four-dimensional Gaussian vector independent of $Pu$.
Conditionally on $Pu$, the projection of $H$ onto $(Pu)^\perp$ is a standard three-dimensional Gaussian.
On the event $\norm P\le B_M$, a projected directional angle at most $a\delta$ implies
\[
 \big|\operatorname{proj}_{(Pu)^\perp}H\big|
 \le \frac{C_{M,a}\delta}{\sin\theta}.
\]
The three-dimensional Gaussian small-ball bound consequently gives
\begin{equation}\label{eq:collision-probability}
 \Pr\{\norm P\le B_M,\ \ang(Pu,Pv)\le a\delta\}
 \le C_{M,a}\min\{1,(\delta/\theta)^3\}.
\end{equation}
The norm constraint has not been treated as an independent event; the event $Pu=0$ has probability zero.

For each original direction, decompose into angular annuli $\theta\asymp2^j\delta$.
By\eqref{eq:angular-cap}, the expected number of collisions is at most
\begin{equation}\label{eq:collision-degree}
 C_{M,s,a}A\sum_{0\le j\le C+\log_2(1/\delta)}2^{j(s-3)}
 \le Q:=C_{M,s,a}A\delta^{3-s}\Ldelta.
\end{equation}
The contribution of projective angles larger than a fixed constant is controlled in the same way,
using a fixed finite cover of the direction space and $m\le C_MA\delta^{-s}$.
The logarithmic loss is present when $s=3$; for $s>3$ the slightly weaker uniform expression above suffices.

The expected number of short directions is less than $m/128$.
The expected number of collision-graph edges counted on the bounded-norm event is at most $CmQ$.
Applying Markov's inequality to these counts and combining the failure probabilities,
we can choose a single $P$ with bounded norm, at least $3m/4$ nonshort directions,
and at most $CmQ$ collision-graph edges among these directions.
The graph vertices are original labels, with an edge whenever the projected angle is at most $a\delta$.
The independent-set bound
\[
 \alpha(G)\ge\frac{|V(G)|^2}{|V(G)|+2|E(G)|}
\]
yields\eqref{eq:projection-count}, together with\eqref{eq:nonshort} and\eqref{eq:projected-separated}.
\end{proof}

\begin{remark}
The lemma does not assume that the projected directional distribution automatically inherits a Frostman upper bound.
Only after selecting an independent set in the collision graph do the separated directions in $\R^4$
have angular upper bounds at all scales, by volume packing in three-dimensional projective direction space.
The lemma also supplies no additional dimension along the projection kernel.
\end{remark}

\begin{lemma}[Average density, a common chart, and the visited union]\label{lem:marks}
Suppose that the original directions are $\delta$-separated and satisfy\eqref{eq:angular-cap},
and that\eqref{eq:single-line-upper} and\eqref{eq:mean-density} hold.
One can first select at least $c_{M,R}\lambda m$ original lines and then apply Lemma~\ref{lem:projection}
to obtain a subfamily $\T'$ in a common bounded four-dimensional graph chart.
On a grid of size $\rho\asymp\delta$, it satisfies
\begin{align}
 m':=\card{\T'}&\ge
 \frac{c_{M,R,s}\lambda m}{A\delta^{3-s}\Ldelta},\label{eq:heavy-output-count}\\
 \card{J'_T}&\ge c_{M,R}\lambda\delta^{-1}
       &&(T\in\T'),\label{eq:individual-output-marks}\\
 \card{E'}&\le C_{M,R,s}\card{E}.\label{eq:output-union}
\end{align}
Every output mark comes from its own original label; extending supporting lines adds no shading.
\end{lemma}

\begin{proof}
First remove the light lines with $\card{Y(T)}<(\lambda/2)\delta^{-1}$.
Their total mass is at most $\lambda m\delta^{-1}/2\le S/2$.
Together with\eqref{eq:single-line-upper}, this shows that at least $c_{M,R}\lambda m$ heavy lines remain.
Each retained line has at least $(\lambda/2)\delta^{-1}$ of its own original marks.
The directional upper bound\eqref{eq:angular-cap} persists on subfamilies without renormalization.
Apply Lemma~\ref{lem:projection} to the heavy family.

Move each original shaded cell to the nearest parameter point on its original line,
and remove a fixed multiplicity in the original line parameter.
Since $|Pv_T|\ge c_M$, every projected line has a directional coordinate whose absolute value is bounded below by a fixed positive constant.
Partition into finitely many charts according to a largest component and its sign, and choose a chart containing the most lines.
Use that coordinate as the common time variable. The other slopes and intercepts then have fixed bounds.
Only a fixed number of originally $\delta$-separated parameter marks can fall in a new time cell of length $\rho=C\delta$.
Thus each line retains the number of marks in\eqref{eq:individual-output-marks}.

All marks lie in a fixed bounded region. Divide the new time axis into finitely many unit windows,
let each line choose the window containing the most of its own marks,
and then pigeonhole the lines by window.
The number of marks per line and the number of lines each lose only a fixed proportion.
Translate the common window to $[0,1]$ and extend supporting lines if necessary, retaining only the existing marks within the window.
This also handles shadings concentrated near segment endpoints.
Choosing $a$ in the projection lemma to be a sufficiently large fixed constant ensures that the final slopes are separated at scale $\rho$.

Finally, the bounded-norm map $P$ sends each original $\delta$-cell into a set of diameter $O_M(\delta)$,
which is covered by a fixed number of four-dimensional $\rho$-cells.
Moving to nearest line points, changing charts, and regrouping the time grid increase the cell count only by fixed factors.
This gives\eqref{eq:output-union} and\eqref{eq:heavy-output-count}.
Marks from different original lines have not been spliced together on a new line.
\end{proof}

\section{The unsaturated finite estimate and the six-dimensional recurrence}\label{sec:seeds}

\begin{proposition}[A finite estimate retaining the original number of lines]\label{prop:finite}
Assume Input~\ref{input:D4}. Fix $M\ge4$, $R$, a real $s\ge3$, and $\varepsilon_D>0$.
There exist $\eta_0>0$ and a scale threshold such that, whenever $m$ original lines with separated directions satisfy
\eqref{eq:angular-cap},\eqref{eq:mean-density}, and $\lambda\ge\delta^{\eta_0}$, one has
\begin{equation}\label{eq:finite-sharp-budget}
 \card{E}\ge
 c_{M,R,s,\varepsilon_D}
 \frac{\lambda^2m}{A\Ldelta}\,
 \delta^{s-4+\varepsilon_D}.
\end{equation}
The number $m$ need not saturate the directional count.
\end{proposition}

\begin{proof}
Obtain a four-dimensional subfamily from Lemma~\ref{lem:marks}.
Its output slopes are separated, so the total unit weight in each slope cell is bounded by a fixed constant.
First choose $\eta_D$ from the output loss in Input~\ref{input:D4}, then choose
$\eta_0<\eta_D/2$, and finally take $\delta$ sufficiently small.
On the comparable grid $\rho\asymp\delta$, the marks on each output line meet the threshold of the four-dimensional input;
all fixed constants are absorbed by this strict margin.
Input~\eqref{eq:D4} therefore yields
\[
 \card{E}\ge c\card{E'}
 \ge c\delta^{\varepsilon_D}\sum_{T\in\T'}\card{J'_T}
 \ge c\delta^{\varepsilon_D}m'\lambda\delta^{-1}.
\]
Substituting\eqref{eq:heavy-output-count} gives\eqref{eq:finite-sharp-budget}.
The two factors of $\lambda$ come respectively from the number of heavy lines and the original marks on each output line;
neither factor can be omitted.
\end{proof}

\begin{corollary}[An explicit small-loss form]\label{cor:uniform}
Fix an integer $M\ge4$, a spatial bound $R$, and a real $s\ge3$. For every $\varepsilon>0$,
there exist $\eta>0$, $C$, and $\delta_0>0$ such that
\begin{equation}\label{eq:tolerance-assumptions}
 A\le\delta^{-\eta},\qquad \lambda\ge\delta^\eta
\end{equation}
implies
\begin{equation}\label{eq:unsaturated-finite}
 \card{E}\ge C^{-1}m\delta^{s-4+\varepsilon}.
\end{equation}
These constants are independent of $m$.
\end{corollary}

\begin{proof}
First take $\varepsilon_D<\varepsilon/4$, then choose
$0<\eta<\min\{\eta_D/2,\varepsilon/16\}$, decreasing it further if necessary.
In\eqref{eq:finite-sharp-budget}, $\lambda^2/A\ge\delta^{3\eta}$.
Finally, take $\delta$ sufficiently small and use the remaining positive exponent margin to absorb $\Ldelta$, the fixed grids, and the constants.
This gives $\forall\varepsilon\,\exists\eta$, without an additional linear modulus.
\end{proof}

\begin{corollary}[Seeds for every fixed real directional parameter]\label{cor:seed}
Under the same named inputs, for every fixed real $n\ge4$ and every $0<\zeta<1$,
$\Kest(n,4-\zeta)$ holds in every fixed ambient space.
\end{corollary}

\begin{proof}
Take $s=n-1\ge3$.
Given a target loss $\tau>0$, first choose a sufficiently small output loss in Corollary~\ref{cor:uniform},
then choose the input tolerance $\eta$ to be smaller than a sufficiently small fixed fraction of $\tau$ as well.
If the saturated count lower bound $m\ge\delta^{-s+\eta}$ also holds, then
\[
 \card{E}\ge c\delta^{-4+\varepsilon+\eta}
            \ge c\delta^{-4+\tau}.
\]
This is an explicit uniform finite sufficient estimate.
Under the arbitrarily small power conventions of Katz--Tao, first fix $\zeta$,
choose the output loss above to be strictly smaller than $\zeta$ by a fixed factor,
and then place all underlying input losses below the corresponding threshold.
This gives the required strictly subcritical $\Kest(n,4-\zeta)$.
When an intermediate family is unsaturated,\eqref{eq:unsaturated-finite} still gives
\[
 \card{E}\ge c\,mN^{4-(n-1)-\varepsilon},
\]
retaining the original number-of-lines factor required by\eqref{eq:KT-unsaturated}.
The ambient dimension remains fixed throughout. If an ambient space has no direction family of the saturated size,
the corresponding statement is vacuous.
\end{proof}

\begin{proposition}[The actual six-dimensional parameter substitution]\label{prop:six}
For every $0<\zeta<1$,
\begin{equation}\label{eq:six-finite}
 \Kest\!\left(6,\frac{17}{4}-\frac{\zeta}{4}\right)
\end{equation}
holds.
\end{proposition}

\begin{proof}
Corollary~\ref{cor:seed} gives
\[
 \Kest(6,4-\zeta),\qquad \Kest(5-\zeta,4-\zeta).
\]
The second real directional parameter is at least $4$, and its angular exponent is $4-\zeta\ge3$,
so it is indeed within the range of the same finite projection estimate.
Taking $n=6$ and $d=d'=4-\zeta$ in Input~\ref{input:KT} yields
\[
 \frac{2n+1+d'}4=\frac{13+4-\zeta}{4}
                  =\frac{17}{4}-\frac\zeta4.
\]
Both seed applications and the recurrence use fixed parameters:
first fix the final output loss, then choose the input tolerances for the finitely many applications,
and finally take the scale sufficiently small.
We invoke neither an unproved critical uniform modulus nor a free monotonicity principle in directional dimension.
\end{proof}

\section{The Hausdorff conclusion for arbitrary sets}\label{sec:hausdorff}

\begin{lemma}[From a finite estimate to arbitrary sets]\label{lem:hausdorff}
Suppose that the finite estimate $\Kest(q,d)$ holds for an integer $q\ge2$ under the arbitrarily small power conventions above.
If $K\subset\R^q$ contains a unit line segment in every direction, then $\HH K\ge d$.
Neither $K$ nor a direction-to-line selection map is required to be measurable.
\end{lemma}

\begin{proof}
Suppose, to the contrary, that $\HH K<t<d$ for some $t$.
First choose an open chart of near-vertical directions parameterized by slopes $v\in\R^{q-1}$.
Every candidate unit segment has a time-coordinate projection of fixed positive length,
and hence contains a time interval with rational endpoints. Its intercept lies in some integer-bounded box.
There are only countably many combinations of these time intervals and intercept boxes,
and their sets of admissible slopes cover the entire direction chart.
By countable subadditivity of Lebesgue outer measure, one combination has an admissible slope set $S_0$
of positive finite outer measure $a>0$.
Apply a fixed affine normalization so that the common time interval is $[0,1]$
and all relevant trajectories lie in a fixed bounded region.
At this point no witness line has been selected for every slope.

Since $\HH K<t$, cover $K$ by balls $B_i$ with radii satisfying
\begin{equation}\label{eq:hausdorff-cover}
 \sum_i r_i^t\le1,\qquad\sup_i r_i\longrightarrow0.
\end{equation}
The same type of cover may be taken for the fixed normalized region; the fixed affine transformation changes only constants.
Let $U_k$ be the union of the balls with $2^{-k-1}<r_i\le2^{-k}$.
There are at most $C2^{kt}$ such balls.

Choose $c_0>0$ so that $\sum_{k\ge1}c_0k^{-2}<1$.
Every successful witness segment $\ell([0,1])\subset K$ satisfies
\[
 \sum_k\big|\{u\in[0,1]:\ell(u)\in U_k\}\big|\ge1.
\]
Thus, for some $k$, its time length in $U_k$ is at least $c_0k^{-2}$.
Let $S_k$ be the set of slopes admitting such a bounded witness line.
The set $S_k$ need not be measurable; nevertheless, $S_0\subset\bigcup_kS_k$.
Countable subadditivity of outer measure again shows that, for one common $k$,
\begin{equation}\label{eq:outer-direction}
 |S_k|^*\ge c_1 a k^{-2},
\end{equation}
where $c_1>0$ is a fixed constant.
Thus the same scale has both an outer directional-mass lower bound and a time-length lower bound on each witness line.

Set $\delta=2^{-k}$, and take a maximal finite $C\delta$-separated set of slopes in the bounded set $S_k$.
Its $C\delta$-balls cover $S_k$, so\eqref{eq:outer-direction} gives
\begin{equation}\label{eq:finite-witness-count}
 m\ge c a k^{-2}\delta^{-(q-1)}.
\end{equation}
A fixed choice of $C$ ensures that the actual unit directions are also separated at scale $\delta$.
Only now choose, for each of these finitely many slopes, a successful witness line.
Packing gives the directional upper bound at every scale:
\[
 \card{\{T:\ang(v_T,u)\le\theta\}}
 \le C_q(\theta/\delta)^{q-1}.
\]
Each line has at least $c\delta^{-1}k^{-2}$ nonempty time cells in $U_k$.
Choose an actual point of the line in each such cell and move it to the nearest grid-cell centre.
Each original covering ball, whose radius is at most $\delta$, contributes only a fixed number of visited cells. Consequently,
\begin{equation}\label{eq:cover-upper}
 \card{E_k}\le C\delta^{-t}.
\end{equation}
The nearest cell centre is within distance $C\delta$ of the original line,
as allowed by the discrete supporting-line convention.
If the convention uses unit arclength segments, split a graph segment of fixed length into finitely many unit-scale windows
and retain the richest window on each line. This loses only a fixed proportion and adds no marks.

By\eqref{eq:finite-witness-count} and the marks on each line,
directional saturation and shading saturation lose only polylogarithmic factors.
First fix $0<\varepsilon<d-t$, then choose the input tolerance for the finite estimate, and finally take $k$ sufficiently large.
All factors $k^{-2}$ and all fixed constants fall within the arbitrarily small power thresholds, so the finite lower bound gives
\begin{equation}\label{eq:cover-lower}
 \card{E_k}\ge c\delta^{-d+\varepsilon}.
\end{equation}
As the largest covering radius tends to zero, the selected $k$ must tend to infinity.
Equations~\eqref{eq:cover-upper} and~\eqref{eq:cover-lower} then contradict each other.
Throughout the argument, ordinary time length is measured only on preimages of open unions of balls along individual lines.
We never integrate over the possibly nonmeasurable set $K$ or use a measurable direction-to-line selection.
\end{proof}

\begin{proof}[Proof of Theorem~\ref{thm:main}]
For every $0<\zeta<1$, Proposition~\ref{prop:six} and Lemma~\ref{lem:hausdorff} give
\[
 \HH K\ge\frac{17}{4}-\frac\zeta4.
\]
This holds for every such fixed $\zeta$. Letting $\zeta\downarrow0$ proves\eqref{eq:main}.
We first establish the finite estimate and the Hausdorff conclusion for each fixed parameter and only then take the parameter limit;
a covering-number lower bound at one fixed scale has not been identified directly with a Hausdorff dimension bound.
\end{proof}

\section{Scope, attribution, and limitations}\label{sec:scope}

The assertion for arbitrary ordinary six-dimensional Kakeya sets is a consequence of the explicitly stated four-dimensional finite input
and the named Katz--Tao recurrence.
For a complete verification, the reader must also examine the original proof of the four-dimensional finite interface and the advanced inputs it cites.
Author attribution, archival hashes, numerical experiments, and internal review labels are not proofs of correctness.

Three levels must be distinguished: the full-dimensional theorem for four-dimensional sets;
the four-dimensional finite directional-mass estimate for labelled original marks;
and its higher-dimensional consequences through the recurrence.
This paper uses the second level as its actual input.
Its output applies uniformly to each new finite family; the families need not arise from one common set across scales.
However, it imposes a long-mark threshold on each line and therefore does not assert a sharp Kakeya maximal estimate for arbitrarily small shading densities.

Directional selection retains the original angular upper bound as an unnormalized subfamily bound.
Passing to a normalized probability law requires division by the retained mass.
The average shading lower bound first undergoes heavy-line selection before entering the per-line input,
so\eqref{eq:finite-sharp-budget} indeed has two factors of $\lambda$.
Projected directions become separated only after selection in the collision graph;
projected Frostman control is not inherited for free.
Arbitrary fixed ambient spaces and real angular parameters are genuine requirements of the recurrence, rather than merely changes in terminology.

Relative to the ordinary six-dimensional Hausdorff baseline explicitly listed in\cite{KT,HRZ,GLX},
equation~\eqref{eq:comparison} gives a fixed positive improvement in the exponent.
The four-dimensional manuscript\cite[Section 11.2]{OAI4} explicitly records the lower bound $4$ obtained by projection in higher ambient dimensions.
The bounded comparison through the stated verification cutoff did not identify an explicit report of the same $17/4$
or a concrete result of the same type that dominates it.
Finite searches, however, establish neither absolute priority nor current optimality, and do not exclude implicit consequences elsewhere.

We do not propose a new Katz--Tao recurrence or claim independent credit for the four-dimensional manuscript or its advanced inputs.
This paper does not settle full Hausdorff dimension in five dimensions or in general dimension, nor does it construct a genuine counterexample.
The manuscript is intended for further external mathematical review.
Internal cross-review is distinguished strictly from external peer review.

\appendix\section{The source and traceable proof of the four-dimensional finite interface}\label{app:inputs}

\subsection{The initial finite-scale consequence}
The four-dimensional directional-mass input starts from the following proof-level finite statement, archived as SRC-001.
Fix a bounded four-dimensional graph chart. For every $\sigma>0$, there exist $\eta_{\rm src}>0$ and $\delta_0>0$
such that a finite family with $\delta$-separated directions, at least $\delta^{-3+\eta_{\rm src}}$ lines,
and at least $\delta^{-1+\eta_{\rm src}}$ marks on each line satisfies
\begin{equation}\label{eq:source-finite}
 \card{E}>\delta^{-4+\sigma}
\end{equation}
(fixed constants may be absorbed by strict exponent margins).
This statement uses the finite initialization of\cite{OAI4} and all subsequent model-exclusion cases;
it is not an abstract compactness consequence of the Hausdorff dimension conclusion.
The complete assembly and finite conclusion are recorded in\cite{ArchiveSource}.
We do not reproduce the four-dimensional proof here.

The archived verification includes independent smoothing of intercepts and slopes at scale $\delta$;
a subpower upper bound for the three-dimensional slope density from directional separation and the number of lines;
conversion of the actual mark intervals to selected durations; control of finite configuration complexity;
attainment of a positive terminal scale while retaining the original unnormalized mass;
and finally all zero, positive finite, and infinite narrowness cases of the original proof.
Checking only the initial reduction would not certify\eqref{eq:source-finite}.

\subsection{From a dense statement to sparse and weighted forms}
To explain why Input~\ref{input:D4} does not additionally assume a saturated directional count,
we recall the finite reduction in\cite{ArchiveDIR}.
First consider $m$ $\delta$-separated slopes. Their $\delta$-neighbourhood has volume $\asymp m\delta^3$
in three-dimensional slope space.
Using $L\le C\delta^{-3}/m$ translations in a fixed bounded range,
the expected covered volume for independent random translations gives a choice whose union has fixed positive volume.
More explicitly, each point in a fixed target slope box is covered by one translation with probability at least $c m\delta^3$.
With the stated number of translations, the expected covered volume in that box has a fixed positive lower bound.
Thus the translated labelled slopes contain $\asymp\delta^{-3}$ directions separated at a fixed multiple of scale $\delta$.

Each slope translation $v\mapsto v+a$ corresponds to the space-time shear
\[
 (t,y)\mapsto(t,y+ta).
\]
The shear retains each copied line's own time marks.
Each copy visits at most $C\card E$ cells, and the union of all copies visits at most $CL\card E$ cells.
After regrouping a grid of a fixed comparable size and taking the original threshold to be a sufficiently small fixed fraction of $\eta_{\rm src}$,
we can apply\eqref{eq:source-finite} to the amplified family and obtain
\begin{equation}\label{eq:sparse-source}
 \card E\ge c m\delta^{-1+\sigma}.
\end{equation}
All shear ranges and coordinate bounds are fixed; no small-power geometric loss is introduced.

For a general weighted family, colour the slope cells with a fixed finite number of colours so that
any representatives from distinct cells of one colour are separated at a fixed multiple of scale $\delta$.
Choose an original label in each nonempty cell and retain its own marks; the representatives' visited union lies within the original union.
By\eqref{eq:sparse-source}, the total number of nonempty cells is at most $C\delta^{1-\sigma}\card E$.
On the other hand, every cell has total weight at most $\alpha$, and each line has at most $C\delta^{-1}$ marks. Hence
\[
 \sum_i\omega_i\card{J_i}
 \le C\delta^{-1}\alpha\cdot
       (\text{number of nonempty slope cells})
 \le C\alpha\delta^{-\sigma}\card E.
\]
Taking $\sigma<\varepsilon_D$, decreasing the input threshold appropriately, and absorbing fixed constants gives\eqref{eq:D4}.
The reduction does not stratify by weight size, so it incurs no loss depending on the smallest weight or the range of weights.

\subsection{Named advanced inputs used in the source proof}
The archived assembly of the four-dimensional internal proof\cite{ArchiveSource} uses the following precise versions as advanced inputs:
\begin{enumerate}
 \item Guth--Wang--Zahl\cite[Theorem 1.1]{GWZ}: the convex-clustering set estimate for distinct tubes in three dimensions,
  with its original average shading density and clustering hypotheses; it is not unconditionally replaced by a maximal-function theorem.
 \item The restricted planar dot-product input of Wang--Zahl\cite[Definition 4.2, Theorem 5.2]{WZ}.
 \item The endpoint multilinear Kakeya input of Carbery--Valdimarsson\cite[Theorem 1]{CV}.
 \item The single-scale Fourier-decay input for multiplicative convolutions of Orponen--de Saxc\'e--Shmerkin\cite[Theorem 1.5]{OSS}.
\end{enumerate}
The proof also uses combinatorial estimates with a fixed number of iterations, real-algebraic and semialgebraic decompositions of fixed format,
the area formula, and standard measure-theoretic tools.
The complete proofs of these advanced external theorems are not reproved from scratch here,
and the formalization reported by the authors has not been rerun locally in this project.
The four-dimensional input, these advanced results, and the subsequent Katz--Tao recurrence are kept distinct in this order to avoid circular dependencies.

\subsection{The fixed-version research archive}
The traceable basis is the research-repository commit of October 19, 2026:
\texttt{774c3614ec3e810c127c875a9f5abb60748be848}.
The table locates the proofs and interfaces cited in this paper.
Archive identifiers locate documents; they do not replace mathematical statements.

\begin{longtable}{@{}p{0.16\textwidth}p{0.76\textwidth}@{}}
 \toprule
 Archive & Content and location \\
 \midrule
 \endhead
 RED-001 & R0038: complete source reduction, finite initialization, and assembly of the internal proof; Section 20 gives the outer-measure conclusion for arbitrary sets.\\
 DIR-003 & R0009: sparse-direction shear amplification and weighted slope-cell reduction.\\
 FSD-001 & R0029: finite Gaussian projection, original marks and common time, arbitrary fixed ambient spaces, and the unsaturated estimate.\\
 DIM-001 & R0109: six-dimensional outcome report, including the complete inputs, parameter chain, and primary-source comparison of the same type.\\
 Independent review & Internal project cross-reviews of the R0029 projection and recurrence interfaces, R0038 source assembly and arbitrary-set argument, R0080 retained core, and R0109 outcome scope.\\
 \bottomrule
\end{longtable}
This paper details the projection, finite estimate, and outer-measure argument;
the four-dimensional source estimate and the recurrence remain explicitly cited inputs.
It is therefore an exposition with transparent dependencies, and the omission of the source proof does not imply that every cited result has been completely reproved here.

\begin{thebibliography}{99}
\bibitem{OAI4}
OpenAI.
\emph{Every Four-Dimensional Kakeya Set Has Full Hausdorff Dimension}.
Preprint of September 24, 2026, 175 pp.; fixed repository version:
\href{https://github.com/openai/math/tree/1581d4cbdd69a0faf34b2089614fe2a814316a6a/preprints/Every-four-dimensional-Kakeya-set-has-full-Hausdorff-dimension-September-24-2026}{openai/math, commit 1581d4cb}.

\bibitem{KT}
Nets Hawk Katz and Terence Tao.
\emph{New bounds on Kakeya problems}.
Journal d'Analyse Math\'ematique \textbf{87} (2002), 231--263.
Preprint version used here: \href{https://arxiv.org/abs/math/0102135v1}{arXiv:math/0102135v1}, 2001;
Sections 5--6, especially Theorem 6.2.

\bibitem{HRZ}
Jonathan Hickman, Keith M. Rogers and Ruixiang Zhang.
\emph{Improved bounds for the Kakeya maximal conjecture in higher dimensions}.
\href{https://arxiv.org/abs/1908.05589v1}{arXiv:1908.05589v1}, 2019; especially Figure 5.

\bibitem{GLX}
Chuanwei Gao, Diankun Liu and Yakun Xi.
\emph{Curved Kakeya sets and Nikodym problems on manifolds}.
\href{https://arxiv.org/abs/2503.11574v1}{arXiv:2503.11574v1}, 2025;
especially the Euclidean comparison in Table 1.

\bibitem{GWZ}
Larry Guth, Hong Wang and Joshua Zahl.
\emph{A streamlined proof of the Kakeya set conjecture in $\R^3$}.
\href{https://arxiv.org/abs/2601.14411v2}{arXiv:2601.14411v2}, 2026; Theorem 1.1.

\bibitem{WZ}
Hong Wang and Joshua Zahl.
\emph{Sticky Kakeya sets and the sticky Kakeya conjecture}.
\href{https://arxiv.org/abs/2210.09581v2}{arXiv:2210.09581v2}; Definition 4.2 and Theorem 5.2.

\bibitem{CV}
Anthony Carbery and Stef\'an Ingi Valdimarsson.
\emph{The endpoint multilinear Kakeya theorem via the Borsuk--Ulam theorem}.
\href{https://arxiv.org/abs/1205.6371v1}{arXiv:1205.6371v1}, 2012; Theorem 1.

\bibitem{OSS}
Tuomas Orponen, Nicolas de Saxc\'e and Pablo Shmerkin.
\emph{On the Fourier decay of multiplicative convolutions}.
\href{https://arxiv.org/abs/2309.03068v2}{arXiv:2309.03068v2}; Theorem 1.5.

\bibitem{ArchiveSource}
Research archive of this project.
\emph{R0038: Arbitrary-set reduction, assembly of the original proof, and finite/Hausdorff conclusions}.
October 17, 2026; fixed-version
\href{https://github.com/lixiang90/kakeya_research/blob/774c3614ec3e810c127c875a9f5abb60748be848/notes/10_source_audit/source_assembly/R0038_reduction_source_assembly_and_finite_readout_full_proof.md}{complete RED-001 proof}.
Internally cross-reviewed; not externally peer-reviewed.

\bibitem{ArchiveDIR}
Research archive of this project.
\emph{R0009: Sparse-direction amplification and weighted directional reduction}.
Fixed-version
\href{https://github.com/lixiang90/kakeya_research/blob/774c3614ec3e810c127c875a9f5abb60748be848/notes/40_multiscale/plank_inputs/R0009_sparse_direction_amplification.md}{sparse-direction proof} and
\href{https://github.com/lixiang90/kakeya_research/blob/774c3614ec3e810c127c875a9f5abb60748be848/notes/40_multiscale/plank_inputs/R0009_weighted_direction_reduction.md}{weighted reduction}.

\bibitem{ArchiveFSD}
Research archive of this project.
\emph{R0029: Finite projection in arbitrary ambient spaces, original marks, and six-dimensional seeds}.
Fixed-version
\href{https://github.com/lixiang90/kakeya_research/blob/774c3614ec3e810c127c875a9f5abb60748be848/notes/20_dimension_lifting/finite_seed_transfer/R0029_projection_and_prior_audit.md}{finite projection proof} and
\href{https://github.com/lixiang90/kakeya_research/blob/774c3614ec3e810c127c875a9f5abb60748be848/notes/20_dimension_lifting/finite_seed_transfer/R0029_conditional_seed_and_six_dimensional_bound.md}{seed and recurrence integration}.
The original compact-set formulation in R0029 was extended by the arbitrary-set conclusion in R0038.

\bibitem{ArchiveDIM}
Research archive of this project.
\emph{DIM-001: A Hausdorff lower bound of 17/4 for arbitrary six-dimensional Kakeya sets}.
October 19, 2026;
\href{https://github.com/lixiang90/kakeya_research/blob/774c3614ec3e810c127c875a9f5abb60748be848/results/reviewed/DIM-001_six_dimensional_Hausdorff_research_report.md}{fixed-version outcome report and review references}.
\end{thebibliography}

\end{document}
